% Generated draft: requires mathematical and source review.
\documentclass[11pt]{article}
\usepackage[margin=25mm]{geometry}
\usepackage{fontspec}
\setmainfont{Latin Modern Roman}
\usepackage{amsmath,amssymb,mathtools,unicode-math}
\setmathfont{Latin Modern Math}
\usepackage{enumitem,needspace,xurl,hyperref}
\setlength{\parindent}{0pt}
\setlength{\parskip}{5pt}
\setlength{\emergencystretch}{3em}
\newcommand{\N}{\mathbb{N}}
\newcommand{\Z}{\mathbb{Z}}
\newcommand{\Q}{\mathbb{Q}}
\newcommand{\R}{\mathbb{R}}
\newcommand{\C}{\mathbb{C}}
\newcommand{\F}{\mathbb{F}}
\newcommand{\A}{\mathbb{A}}
\newcommand{\PP}{\mathbb{P}}
\newcommand{\E}{\mathbb{E}}
\newcommand{\eps}{\varepsilon}
\newcommand{\dd}{\,\mathrm{d}}
\begin{document}
\section*{UnsolvedMath ID: 30005226}
\textbf{Source problem code:} catalogue entry
\subsection{Definitions and mathematical statement}
A partition of a positive integer \(m\) is a finite weakly decreasing sequence of positive integers with sum \(m\); write \(\rho\vdash m\). The partition \((m-1,1)\) is defined for \(m\geq 2\). These conventions clarify the partition terminology without supplying the missing parameter restrictions of the source problem.

For context, if \(\alpha,\beta,\gamma\) are partitions and \(s_\alpha\) denotes the corresponding Schur function, a plethysm multiplicity is a coefficient
\[
[s_\gamma]\bigl(s_\alpha[s_\beta]\bigr),
\]
where brackets inside the argument denote plethysm and \([s_\gamma]\) denotes coefficient extraction in the Schur basis. The supplied record does not identify the Schur functions or indexing partitions that define its quantity \(A_{\ell,m}(\nu,\lambda)\), so no identification with a particular coefficient is asserted here.

\textbf{Source question, with unresolved qualifications.} Determine the limit of
\[
A_{\ell,m}(\nu,\lambda),
\]
and ask whether the associated stability results extend from the partition \((m-1,1)\) to arbitrary partitions \(\rho\vdash m\), through arbitrary maximal tableaux. The workshop question about this extension is attributed to Wildon. The original excerpt describes a possible connection between Wildon's conjecture concerning the form of an expression involving \(\nu\) and plethystic semistandard tableaux of maximal weight, defined in its reference [5]. It gives a corrupted tableau-shape expression beginning with \((m-1,1)\) and involving \(1\) and \(j\); its precise indexing cannot be recovered from the supplied text.

The record supplies neither a definition of \(A_{\ell,m}(\nu,\lambda)\), domains and compatibility conditions for \(\ell,m,\nu,\lambda\), nor the variable and regime of the proposed limit. It also omits the definition of maximality, the relevant tableau shapes and weights, and the precise stability assertion to be generalized. Consequently, this excerpt does not determine a fully quantified limit problem or a formal generalization conjecture. In particular, no limiting variable, formula for the limit, or universal positivity assertion is imposed in this draft.

The supplied literature assessment labels the problem partially solved: it reports that eventual constancy of the displayed plethysm multiplicities is proved, while determining the limit, strict positivity, and the extension to arbitrary maximal tableaux remain open. This is supplied dated curation, not an independently verified mathematical or literature-status assertion; the missing definitions prevent specifying its exact scope.
\subsection{Short English statement}
The source asks for the limiting value of certain plethysm multiplicities and whether stability associated with \((m-1,1)\) extends to arbitrary partitions of \(m\) and maximal tableaux. The supplied assessment reports eventual constancy but leaves the limit, strict positivity, and the broader extension open. The multiplicities, limiting regime, and tableau conditions are missing, so a precise quantified formulation requires further source material.
\subsection{Sources}
\begin{itemize}
\item Supplied local UnsolvedMath registry record, permanent ID \(30005226\), title \emph{Stable Limits of Plethysm Multiplicities}, exact source problem code \texttt{catalogue entry}. The record gives proposed year 2022, no named proposer in its proposer field, and attributes the workshop generalization question to Wildon in the original excerpt. Its supplied status is \texttt{partially\_solved}; its statement status is \texttt{corrected\_verified}, with the qualification that only formula delimiters, notation, grammar, and typography were repaired while mathematical content and scope were preserved. These are registry annotations, not independent verification.
\item Supplied report citation: \emph{Character Theory and Categorification} (2023), Oberwolfach Reports, DOI 10.4171/owr/2022/39, pp. 2267--2333. \url{https://doi.org/10.4171/owr/2022/39}. The record describes Oberwolfach Reports as documenting mathematical workshops and their open-problem discussions.
\item Supplied literature source: Law--Okitani, \emph{Some stable plethysms}. \url{https://doi.org/10.1090/proc/16556}.
\item Supplied literature source: Paget--Wildon, \emph{two stability theorems}. \url{https://arxiv.org/abs/2505.06104}.
\item Supplied literature source: \emph{Oberwolfach Report 39/2022}. \url{https://ems.press/content/serial-article-files/46975}.
\item Literature assessment date: \texttt{2026-08-22}. The assessment is supplied dated curation derived from the supplied OWR curation, not independently verified. It reports proved eventual constancy and unresolved questions about the limit, strict positivity, and generalization to arbitrary maximal tableaux. The source explicitly calls for independent verification of time-sensitive status and cited claims.
\item The original excerpt invokes reference [5] for plethystic semistandard tableaux, but its bibliographic entry is absent from the supplied record. It also includes a figure caption referring to Young diagrams for expressions written as \(\alpha_{\ell,m}\{j\}\) and \(\alpha_{\ell,m}[j]\); their definitions and the figure are not supplied.
\end{itemize}
\end{document}
